% =====================================================================
%  Section 14 — AI in the Project
% =====================================================================
\makesectioncover{AI in the Project}

\section{AI in the Project}
\label{sec:ai}

\subsection{Why AI at all}

Cairo Public Transport is, at its core, a city-scale data-integration,
segmentation, and spatial-pattern problem. Each of those three problem
classes has a classical method that nearly works and a non-classical
method that closes the last 20\% of the gap. The project picks the
classical method first and adds an AI technique only where the
classical method fails on a concrete, named case.

There are three AI techniques in the project. Every one of them is the
answer to a problem the cleaning notebook ran into and could not
otherwise solve cleanly. We document each as: \emph{What problem it
solves}, \emph{Why classical methods fail there}, \emph{How it works},
\emph{Why it is effective for this problem}, and \emph{What we do not
claim}.

% --------------------------------------------------------------------
\subsection{AI Technique 1 \textbullet\ SBERT multilingual semantic matching}

\textbf{Where it lives.} \fn{Phase2/Phase2\_Scraping\_Cleaning.ipynb}
— Stage 4 of the four-stage integration pipeline (see
\S\ref{sec:phase2-integration}).

\textbf{Problem it solves.} Cross-script (Arabic $\leftrightarrow$ English)
station-name matching. When integrating Wikipedia metro stations
(English-first), Google Maps BRT scrapes (mixed), and Phase 1
informal terminals (Arabic-first), the same physical mawfaq can be
written as:

\begin{itemize}
  \item \ar{عدلي منصور} (Arabic) — \strval{Adly Mansour} (English) — \strval{Adli Mansour} (transliterated)
  \item \ar{ميدان مصر الجديدة} — \strval{Heliopolis Square}
  \item \ar{ملعب القاهرة} — \strval{Cairo Stadium}
\end{itemize}

\textbf{Why classical methods fail here.} RapidFuzz token-sort uses
character-level edit distance; Arabic and English share zero
characters so the score is always near zero. \fn{uroman}
transliteration helps with phonetic equivalents (\ar{المرج} $\to$
\strval{almrj}) but breaks on the cases above where the cross-language
form is \emph{semantic} rather than phonetic
(\strval{Heliopolis Square} is not a phonetic rendering of
\ar{ميدان مصر الجديدة}).

\textbf{How it works.}
\begin{enumerate}
  \item Each station name is encoded by
        \fn{sentence-transformers/paraphrase-multilingual-MiniLM-L12-v2}
        into a \stat{384-dim} embedding vector.
  \item The model was pre-trained on parallel sentence pairs across 50+
        languages including Arabic and English, so semantically
        equivalent sentences land near each other in the embedding
        space regardless of script.
  \item Pairwise cosine similarity is computed for every candidate
        pair surviving Stages 1--3.
  \item Pairs above $\tau = 0.65$ are accepted as matches.
\end{enumerate}

\begin{lstlisting}[caption={Stage-4 SBERT matching, simplified.}]
from sentence_transformers import SentenceTransformer
import numpy as np

model = SentenceTransformer(
    "sentence-transformers/paraphrase-multilingual-MiniLM-L12-v2"
)
emb_p1 = model.encode(phase1_names)   # (n1, 384)
emb_p2 = model.encode(phase2_names)   # (n2, 384)

# Cosine similarity matrix (n1 x n2)
e1 = emb_p1 / np.linalg.norm(emb_p1, axis=1, keepdims=True)
e2 = emb_p2 / np.linalg.norm(emb_p2, axis=1, keepdims=True)
sim = e1 @ e2.T

mask = sim >= 0.65   # threshold
\end{lstlisting}

\textbf{Why effective here.}
\begin{itemize}
  \item Egyptian station names are \emph{bilingual by convention} —
        the same node will appear in both languages across our sources.
        SBERT was made for exactly this.
  \item MiniLM is small (384-dim, $<$120\,M parameters) and fast on CPU
        — at $\sim$50 stations per source, we don't need a larger model.
  \item Stages 1--3 catch $\sim$80\% of pairs cheaply; SBERT recovers
        the cross-script tail without paying its cost on the easy
        cases. The cascade structure is what makes the AI technique
        affordable.
\end{itemize}

\textbf{What we do not claim.} The threshold $\tau = 0.65$ is calibrated
on a 5-pair gold set — that is small. Section~\ref{sec:phase2-integration}
records the honest bound: SBERT matches are \emph{strong candidates},
not ground truth, and every Phase 2 hypothesis has been run with the
SBERT-only matches dropped to confirm robustness.

% --------------------------------------------------------------------
\subsection{AI Technique 2 \textbullet\ K-Means district segmentation}

\textbf{Where it lives.} \fn{Phase2/Phase2\_Analysis\_Hypothesis.ipynb}
— Q24.

\textbf{Problem it solves.} Multidimensional district segmentation.
Single-variable rankings (``rank by population'', ``rank by station
count'') cannot describe the joint structure of demand, density,
growth, and infrastructure. Without a joint segmentation we cannot
size a market.

\textbf{Why a single-variable approach fails.} The product story (Section
\ref{sec:synthesis}) needs to identify districts where coverage is
weak \emph{and} demand is high \emph{and} growth is sustained. Any one
of those filters alone produces a different ranking; the joint set
cannot be reproduced by sorting on any single column.

\textbf{How it works.}
\begin{enumerate}
  \item Five engineered features per district:
        \fn{log(Population\_2023)}, \fn{CAGR \%} (2006$\to$2017),
        stations within 3\,km, informal stops within 3\,km,
        Cairo-governorate indicator.
  \item Z-score standardisation so no feature dominates the distance.
  \item K-Means with \fn{k=4}, \fn{init="k-means++"}, \fn{n\_init=50}.
  \item Choice of $k$ defended by an elbow + silhouette scan
        (silhouette at $k=4$ = \stat{0.412}, the local maximum).
  \item Stability validated via re-running across 5 random seeds and
        computing Adjusted Rand Index — \stat{ARI = 1.00} (every seed
        produces the identical partition).
\end{enumerate}

\textbf{Why effective here.}
\begin{itemize}
  \item Continuous z-scored features map cleanly to Euclidean distance
        — the problem is well-posed for K-Means.
  \item $k=4$ is intuitively interpretable: Formal-Served Core,
        Peripheral Growth, Mixed, Low-Activity Outskirts.
  \item \stat{ARI = 1.00} is unusually strong stability — the
        partition is reproducible for any seed, which makes it
        defensible as the foundation for a market-sizing claim.
\end{itemize}

\textbf{What we do not claim.} K-Means is \emph{descriptive}, not
\emph{causal}. The 18.9M-resident addressable market estimated from
the partition is a sizing claim, not a forecast. The clusters are a
useful prism through which to understand the city, not a prediction of
which districts will or will not adopt a product.

% --------------------------------------------------------------------
\subsection{AI Technique 3 \textbullet\ Moran's I spatial autocorrelation}

\textbf{Where it lives.} \fn{Phase2/Phase2\_Analysis\_Hypothesis.ipynb}
— H1 robustness check.

\textbf{Problem it solves.} Once Kruskal-Wallis says ``dense districts
are systematically under-served'' (H1's main result), the next question
is: \emph{is the under-service spatially concentrated, or randomly
scattered through the city?} This question is qualitative when asked
by eye and quantitative when asked by Moran's~I. The answer materially
changes the planning recommendation.

\textbf{Why a single hypothesis test is not enough.} Kruskal-Wallis
treats districts as independent. If the under-served districts cluster
geographically (e.g. all in north-west Cairo), the planning
intervention is local — open a single corridor and the gap closes. If
they are scattered, no single corridor can close the gap. The two
worlds look identical to Kruskal-Wallis.

\textbf{How it works.}
\begin{enumerate}
  \item Build a spatial weight matrix $W$ with
        \fn{pysal.lib.weights.KNN(centroids, k=4)}, row-standardised.
        We use \kw{KNN-as-contiguity} on Nominatim district centroids
        (S6) because Cairo districts are administrative polygons of
        irregular shape and we want neighbour relationships defined
        by centroid proximity, not shared-boundary.
  \item For each district, $z_i = x_i - \bar{x}$ where $x$ is the
        per-district stations-per-100k.
  \item Compute Moran's I:
        $I = \frac{n}{\sum_{i,j} w_{ij}} \cdot
        \frac{\sum_{i,j} w_{ij} z_i z_j}{\sum_i z_i^2}$.
  \item Run \fn{esda.Moran(values, w, permutations=999)} to get an
        empirical p-value.
\end{enumerate}

\textbf{Why effective here.}
\begin{itemize}
  \item Moran's I is the canonical tool for the question. There is no
        cleaner alternative.
  \item The 999-permutation p-value avoids any normality assumption.
  \item It is \emph{additive evidence}, not a substitute for the main
        Kruskal test — if Moran's I had said ``not clustered'' the
        Kruskal result would still hold, just with different planning
        implications.
\end{itemize}

\textbf{What we do not claim.} The Phase 2 result is
\stat{I = 0.087, z = 1.805, p = 0.0500} — borderline-positive.
$n = 68$ districts is at the lower edge of where Moran's I has
statistical bite; with $n = 200$ districts the same effect would have
a much smaller p-value. We treat the result as \emph{suggestive}:
under-service is geographically concentrated, but not so cleanly that
a single corridor would close the gap.

% --------------------------------------------------------------------
\subsection{When AI was decisive and when it wasn't}

\begin{callout}
\textbf{\color{plotblue}Honest reading of the AI contribution.}\quad
SBERT was \emph{decisive} for cross-script integration — without it,
matching \strval{Adly Mansour} to \ar{عدلي منصور} would have required
manual review of every candidate pair. K-Means is \emph{descriptively
powerful and reproducible} but does not prove a causal claim about
market opportunity. Moran's I is \emph{the right tool for the
geographic-clustering question} but the dataset (n=68) is small, so
the p-value lands on the boundary rather than well inside.

The project does not pretend AI made any of its arguments true. It
made specific technical problems tractable. Every conclusion in
Section~\ref{sec:synthesis} rests on the underlying data — the AI
techniques are part of how we got the data into shape.
\end{callout}

\subsection{Where each AI technique sits in the pipeline}

\begin{callout}
\textbf{\color{plotblue}Pipeline location of each AI technique.}\quad
\begin{itemize}
  \item \kw{Raw scraped sources} (S1, S3--S8) $\to$
        per-source cleaning (KNN imputation, \fn{uroman}
        transliteration, type casting) $\to$
        integration stages 1--3 (KNN spatial, OSM cross-verify,
        RapidFuzz).
  \item \kw{AI 1 — SBERT} sits at integration Stage 4. It is the last
        link in the integration cascade.
  \item Cleaned + integrated layer feeds into Phase 2 analysis.
  \item \kw{AI 2 — K-Means} runs at Q24 to produce the 4-cluster
        district segmentation.
  \item \kw{AI 3 — Moran's I} runs as the spatial robustness check on
        H1.
\end{itemize}
\end{callout}
